\documentclass[12pt,a4paper]{article}
\usepackage[UTF8]{ctex}
\usepackage{geometry}
\usepackage{listings}
\usepackage{xcolor}
\usepackage{hyperref}
\usepackage{amsmath}

\geometry{margin=2.5cm}

\lstset{
    language=Python,
    basicstyle=\ttfamily\small,
    keywordstyle=\color{blue}\bfseries,
    commentstyle=\color{green!60!black},
    stringstyle=\color{red},
    numbers=left,
    numberstyle=\tiny\color{gray},
    stepnumber=1,
    numbersep=5pt,
    frame=single,
    breaklines=true,
    showstringspaces=false,
    tabsize=4
}

\title{Halo 系统代码文档：Workers 模块}
\author{代码分析文档}
\date{\today}

\begin{document}

\maketitle

\tableofcontents
\newpage

\section{概述}

\texttt{workers} 模块实现了实际的模型推理执行层。Halo 支持两种 Worker 实现：
\begin{itemize}
    \item \texttt{vLLMWorker}：基于 vLLM 的高性能推理引擎
    \item \texttt{TransformersWorker}：基于 Transformers 库的推理引擎（支持 KV cache 管理）
\end{itemize}

本文档主要分析 \texttt{vLLMWorker}，它是当前系统的主要实现。

\section{vLLMWorker 类}

\subsection{类定义位置}
\texttt{halo/workers/worker\_v.py}

\subsection{功能概述}

\texttt{vLLMWorker} 是一个独立的进程，绑定到特定的 GPU，负责：
\begin{itemize}
    \item 加载和管理模型
    \item 接收执行任务
    \item 执行批量推理
    \item 返回结果给主进程
\end{itemize}

\subsection{初始化}

\begin{lstlisting}
def __init__(
    self,
    id: int,
    physical_gpu_id: int,
    cmd_queue: "queue.Queue",
    result_queue: "queue.Queue",
) -> None:
    self.id = id
    self.device = f"cuda:{physical_gpu_id}"
    
    # 绑定到特定 GPU
    os.environ["CUDA_VISIBLE_DEVICES"] = str(physical_gpu_id)
    
    # 运行时状态
    self.model_name: Optional[str] = None
    self.llm: Optional[LLM] = None
    self.tokenizer: Optional[AutoTokenizer] = None
    
    self.cmd_queue = cmd_queue
    self.response_queue = result_queue
    
    # vLLM 选项
    self.enforce_eager: bool = True
\end{lstlisting}

\begin{description}
    \item[功能] 初始化 Worker，绑定到指定 GPU
    \item[关键点]
    \begin{itemize}
        \item 通过设置 \texttt{CUDA\_VISIBLE\_DEVICES} 环境变量，确保该进程只能看到指定的 GPU
        \item 模型和 tokenizer 初始为 None，按需加载
        \item \texttt{enforce_eager=True} 禁用 CUDA 图优化（提高兼容性）
    \end{itemize}
\end{description}

\subsection{init\_op 方法}

\subsubsection{功能}

为每个 OP 初始化运行时状态，包括：
\begin{itemize}
    \item 采样参数配置
    \item 模型切换（如果需要）
    \item 系统提示词设置
\end{itemize}

\subsubsection{实现代码}

\begin{lstlisting}
def init_op(self, op: Operator) -> None:
    """
    Initialize per-op runtime state:
      - Sampling parameters
      - Optional system/common messages
      - Switch model/tokenizer if op.model_config.model_name changes
    """
    self.op_name = op.id
    self.is_duplicate = getattr(op, "is_duplicate", False)
    
    cfg: ModelConfig = op.model_config
    self.use_chat_template: bool = bool(getattr(cfg, "use_chat_template", True))
    
    # 构建 vLLM 采样参数
    self.sampling_params = SamplingParams(
        temperature=cfg.temperature,
        top_p=cfg.top_p,
        max_tokens=cfg.max_tokens,
        min_tokens=getattr(cfg, "min_tokens", 0),
    )
    
    # 模型切换逻辑
    model_name = cfg.model_name
    if model_name != self.model_name:
        # 释放之前的模型
        if self.llm is not None:
            try:
                del self.llm
            except Exception:
                pass
            self.llm = None
        
        if self.tokenizer is not None:
            try:
                del self.tokenizer
            except Exception:
                pass
            self.tokenizer = None
        
        torch.cuda.empty_cache()
        
        # 加载新模型
        self.llm = LLM(
            model_name,
            dtype=str(getattr(cfg, "dtype", "bfloat16")),
            quantization=getattr(cfg, "quantization", None),
            max_model_len=getattr(cfg, "max_model_len", None),
            enforce_eager=self.enforce_eager,
        )
        self.tokenizer = AutoTokenizer.from_pretrained(model_name, use_fast=True)
        self.model_name = model_name
    
    # 系统提示词
    self.prefix: str = (getattr(cfg, "system_prompt", None) or "").strip()
    self.common_message: str = (getattr(cfg, "common_message", "") or "").strip()
\end{lstlisting}

\subsubsection{关键逻辑}

\begin{enumerate}
    \item \textbf{模型切换检测}：
    \begin{itemize}
        \item 比较当前模型名称和新 OP 的模型名称
        \item 如果不同，释放旧模型并加载新模型
        \item 清空 CUDA 缓存以释放显存
    \end{itemize}
    
    \item \textbf{采样参数}：
    \begin{itemize}
        \item 从 \texttt{ModelConfig} 中提取参数
        \item 创建 vLLM 的 \texttt{SamplingParams} 对象
    \end{itemize}
    
    \item \textbf{提示词设置}：
    \begin{itemize}
        \item \texttt{prefix}：系统提示词
        \item \texttt{common\_message}：通用消息（可选）
    \end{itemize}
\end{enumerate}

\subsection{execute 方法}

\subsubsection{功能}

执行一个 OP 的批量推理任务，这是 Worker 的核心方法。

\subsubsection{方法签名}

\begin{lstlisting}
@torch.inference_mode()
def execute(self, exe_info: ExecuteInfo) -> Dict[str, Any]:
    """
    Execute a single round for one Operator:
      1) Initialize per-op state
      2) Build batch inputs (with or without chat template)
      3) Call vLLM.generate
      4) Return structured results
    
    Returns:
        {
          "item": List[{"id": int, "output": str, "benchmark": (float, float)}],
          "op_name": str,
          "benchmark": {"init_time": float, "prefill_time": float, "generate_time": float}
        }
    """
\end{lstlisting}

\subsubsection{执行流程}

\begin{lstlisting}
def execute(self, exe_info: ExecuteInfo) -> Dict[str, Any]:
    # 1. 初始化 OP
    init_start = time.perf_counter()
    self.init_op(exe_info.op)
    
    # 2. 创建 Query 对象（用于内部处理）
    queries: List[Query] = [
        Query(qid, prompt) for qid, prompt in zip(exe_info.query_ids, exe_info.prompts)
    ]
    init_time = time.perf_counter() - init_start
    
    # 3. 构建输入
    prefill_start = time.perf_counter()
    messages_batch: List[Any] = []
    
    if self.use_chat_template:
        # 使用 HuggingFace 聊天模板
        sys_text = "\n".join([x for x in [self.prefix, self.common_message] if x]).strip()
        for q in queries:
            messages_batch.append([
                {"role": "system", "content": sys_text},
                {"role": "user", "content": q.prompt},
            ])
        inputs: List[str] = self.tokenizer.apply_chat_template(
            messages_batch,
            tokenize=False,
            add_generation_prompt=False,
        )
    else:
        # 普通字符串提示词
        inputs: List[str] = []
        for q in queries:
            header = "\n".join([x for x in [self.prefix, self.common_message] if x]).strip()
            if header:
                inputs.append(f"{header}\n\n{q.prompt}")
            else:
                inputs.append(q.prompt)
    
    # 4. 推理
    outputs = self.llm.generate(inputs, self.sampling_params)
    
    # 5. 收集结果
    results: List[Dict[str, Any]] = []
    for i, output in enumerate(outputs):
        gen_text = output.outputs[0].text if output.outputs else ""
        # 为了可重现性，前置完整输入
        if isinstance(inputs[i], str):
            full_text = inputs[i] + gen_text
        else:
            full_text = "".join([m.get("content", "") for m in inputs[i]]) + gen_text
        
        results.append({
            "id": queries[i].id,
            "output": full_text,
            "benchmark": (prefill_start, time.perf_counter()),
        })
    
    generate_time = time.perf_counter() - prefill_start
    
    # 6. 构建基准数据
    if self.is_duplicate:
        benchmark = {"init_time": 0.0, "prefill_time": 0.0, "generate_time": 0.0}
    else:
        benchmark = {
            "init_time": init_time,
            "prefill_time": 0.0,
            "generate_time": generate_time
        }
    
    return {
        "item": results,
        "op_name": self.op_name,
        "benchmark": benchmark
    }
\end{lstlisting}

\subsubsection{关键步骤详解}

\begin{enumerate}
    \item \textbf{初始化阶段}：
    \begin{itemize}
        \item 调用 \texttt{init\_op} 初始化 OP 状态
        \item 如果需要，加载或切换模型
        \item 记录初始化时间
    \end{itemize}
    
    \item \textbf{输入构建}：
    \begin{itemize}
        \item 如果使用聊天模板：将系统提示词和用户提示词组合成标准聊天格式
        \item 否则：简单拼接系统提示词和用户提示词
        \item 注意：\texttt{exe\_info.prompts} 已经包含了父 OP 的输出（由 Optimizer 拼接）
    \end{itemize}
    
    \item \textbf{批量推理}：
    \begin{itemize}
        \item 调用 \texttt{llm.generate(inputs, self.sampling\_params)}
        \item vLLM 会自动处理批处理和 KV cache
        \item 返回生成结果列表
    \end{itemize}
    
    \item \textbf{结果处理}：
    \begin{itemize}
        \item 提取每个查询的生成文本
        \item 为了可重现性，将完整输入和生成文本拼接
        \item 记录性能基准（时间戳）
    \end{itemize}
    
    \item \textbf{基准数据}：
    \begin{itemize}
        \item 如果是复制节点（数据并行），不记录基准（避免重复计算）
        \item 否则记录 init\_time 和 generate\_time
        \item 注意：vLLM 的 prefill 和 generate 是融合的，所以 prefill\_time 为 0
    \end{itemize}
\end{enumerate}

\subsection{run 方法}

\subsubsection{功能}

Worker 的主循环，持续监听命令队列并执行任务。

\subsubsection{实现代码}

\begin{lstlisting}
def run(self, debug: bool = True) -> None:
    """
    Command loop consuming tasks from cmd_queue and writing
    structured results to response_queue.
    
    Supported commands:
      - "execute": params = (ExecuteInfo,)
      - "exit":    params = ()
    """
    while True:
        msg = self.cmd_queue.get()  # 阻塞等待命令
        
        # 解析消息格式
        if isinstance(msg, tuple):
            command, params = msg
        elif isinstance(msg, dict):
            command = msg.get("command")
            params = msg.get("params", ())
        else:
            self.response_queue.put({
                "command": "error",
                "result": "Unsupported message format",
                "elapsed_time": 0.0
            })
            continue
        
        # 处理退出命令
        if command == "exit":
            out = self.exit()
            self.response_queue.put({
                "command": "exit",
                "result": out,
                "elapsed_time": 0.0
            })
            break
        
        # 动态调用对应方法
        func = getattr(self, command, None)
        if not callable(func):
            self.response_queue.put({
                "command": "error",
                "result": f"Unknown command: {command}",
                "elapsed_time": 0.0
            })
            continue
        
        # 执行命令并返回结果
        start = time.perf_counter()
        try:
            result = func(*params) if params is not None else func()
            elapsed = time.perf_counter() - start
            self.response_queue.put({
                "command": command,
                "result": result,
                "elapsed_time": elapsed
            })
        except Exception as e:
            if debug:
                raise  # 调试模式下抛出异常
            self.response_queue.put({
                "command": "error",
                "result": repr(e),
                "elapsed_time": 0.0
            })
\end{lstlisting}

\subsubsection{消息格式}

所有消息都遵循统一格式：

\begin{lstlisting}
{
    "command": "<command_name>",
    "result": <result_payload>,
    "elapsed_time": <float>
}
\end{lstlisting}

\subsubsection{支持的命令}

\begin{itemize}
    \item \texttt{"execute"}：执行推理任务
    \begin{itemize}
        \item 参数：\texttt{(ExecuteInfo,)}
        \item 返回：包含结果和基准数据的字典
    \end{itemize}
    
    \item \texttt{"exit"}：退出 Worker
    \begin{itemize}
        \item 参数：\texttt{()}
        \item 返回：退出确认消息
    \end{itemize}
\end{itemize}

\subsection{exit 方法}

\begin{lstlisting}
def exit(self) -> str:
    """Release model/tokenizer resources and clear CUDA cache."""
    try:
        if self.llm is not None:
            del self.llm
    except Exception:
        pass
    try:
        if self.tokenizer is not None:
            del self.tokenizer
    except Exception:
        pass
    self.llm = None
    self.tokenizer = None
    torch.cuda.empty_cache()
    logging.info("vLLMWorker[%s] exited.", self.id)
    return "Worker exited."
\end{lstlisting}

\begin{description}
    \item[功能] 清理资源并退出 Worker
    \item[步骤]
    \begin{enumerate}
        \item 删除模型和 tokenizer 对象
        \item 清空 CUDA 缓存
        \item 记录日志
    \end{enumerate}
\end{description}

\section{TransformersWorker 简介}

\texttt{TransformersWorker} 是另一个 Worker 实现，基于 Transformers 库，主要特点：

\begin{itemize}
    \item 支持 KV cache 的显式管理
    \item 支持 cache 在设备间传输
    \item 支持动态批处理大小调整
    \item 支持 prefill 和 decode 的流水线重叠
\end{itemize}

由于当前系统主要使用 \texttt{vLLMWorker}，本文档不详细展开。

\section{性能优化}

\subsection{vLLM 的优势}

\begin{itemize}
    \item \textbf{连续批处理}：自动管理动态批处理，无需手动调整
    \item \textbf{PagedAttention}：高效的 KV cache 管理
    \item \textbf{量化支持}：支持 AWQ、GPTQ 等量化方法
    \item \textbf{高吞吐}：针对批量推理优化
\end{itemize}

\subsection{模型切换优化}

\begin{itemize}
    \item \textbf{延迟加载}：只在需要时加载模型
    \item \textbf{显存管理}：切换模型时主动释放显存
    \item \textbf{调度器优化}：调度器会尽量将相同模型的 OP 分配到同一设备
\end{itemize}

\section{总结}

\texttt{vLLMWorker} 是 Halo 系统的执行层，主要职责包括：

\begin{itemize}
    \item \textbf{模型管理}：按需加载和切换模型
    \item \textbf{批量推理}：执行批量生成任务
    \item \textbf{结果返回}：将结果格式化并返回给主进程
    \item \textbf{性能监控}：记录执行时间等性能指标
\end{itemize}

通过独立进程和队列通信，Worker 实现了与主进程的解耦，支持真正的并行执行。

\end{document}
